\noindent
\begin{minipage}[t]{0.26\textwidth}
\vspace*{14.8cm}
\begin{flushleft}
\footnotesize
Đừng nhầm lẫn Công thức~5 (trong đó $x$ là \textit{số mũ}) với Quy tắc lũy thừa (trong đó $x$ là \textit{cơ số}):
\vspace{-1mm}
\[
\frac{d}{dx}(x^n) = n x^{n-1}
\]
\end{flushleft}
\end{minipage}%
\hfill
\begin{minipage}[t]{0.70\textwidth}

\noindent\textbf{\textsf{\textcolor{stewartcyan}{VÍ DỤ 8}}}\quad Nếu $f(x) = \sin(\cos(\tan x))$, thì
\begin{align*}
f'(x) &= \cos(\cos(\tan x)) \frac{d}{dx} \cos(\tan x) \\[3pt]
&= \cos(\cos(\tan x)) [-\sin(\tan x)] \frac{d}{dx} (\tan x) \\[3pt]
&= -\cos(\cos(\tan x)) \sin(\tan x) \sec^2 x
\end{align*}
Lưu ý rằng chúng ta đã sử dụng Quy tắc chuỗi hai lần.\hfill{\color{stewartcyan}$\blacksquare$}

\vspace{0.4cm}
\noindent\textbf{\textsf{\textcolor{stewartcyan}{VÍ DỤ 9}}}\quad Tính đạo hàm của $y = e^{\sec 3\theta}$.

\vspace{0.15cm}
\noindent\textbf{\textsf{\textcolor{stewartcyan}{LỜI GIẢI}}}\quad Hàm ngoài là hàm mũ, hàm ở giữa là hàm secant, và hàm trong cùng là hàm nhân ba. Do đó ta có
\begin{align*}
\frac{dy}{d\theta} &= e^{\sec 3\theta} \frac{d}{d\theta}(\sec 3\theta) \\[3pt]
&= e^{\sec 3\theta} \sec 3\theta \tan 3\theta \frac{d}{d\theta}(3\theta) \\[3pt]
&= 3e^{\sec 3\theta} \sec 3\theta \tan 3\theta \tag*{{\color{stewartcyan}$\blacksquare$}}
\end{align*}

\vspace{0.4cm}
\noindent{\color{stewartcyan}\rule{1.15ex}{1.15ex}}\quad\textbf{\textsf{\large Đạo hàm của các hàm số mũ tổng quát}}

\vspace{0.2cm}
Chúng ta có thể sử dụng Quy tắc chuỗi để tính đạo hàm của hàm số mũ với cơ số bất kỳ $b > 0$. Nhớ lại từ Phương trình~1.5.10 rằng ta có thể viết
\[
b^x = e^{(\ln b)x}
\]
và khi đó Quy tắc chuỗi cho ta
\begin{align*}
\frac{d}{dx}(b^x) &= \frac{d}{dx}\left(e^{(\ln b)x}\right) = e^{(\ln b)x} \frac{d}{dx}[(\ln b)x] \\[3pt]
&= e^{(\ln b)x}(\ln b) = b^x \ln b
\end{align*}
bởi vì $\ln b$ là một hằng số. Do đó ta có công thức

\vspace{0.25cm}
\noindent
\parbox{\linewidth}{%
  \makebox[0pt][l]{\raisebox{-0.5\height}{\tcbox[colframe=stewartred,colback=white,boxrule=1pt,sharp corners,top=3pt,bottom=3pt,left=6pt,right=6pt]{\textbf{\textcolor{stewartred}{\normalsize 5}}}}}%
  \hfill
  \raisebox{-0.5\height}{\tcbox[colframe=stewartred,colback=white,boxrule=1pt,sharp corners,top=6pt,bottom=6pt,left=24pt,right=24pt]{%
    $\displaystyle \frac{d}{dx}(b^x) = b^x \ln b$%
  }}%
  \hfill\null
}
\vspace{0.35cm}

\noindent\textbf{\textsf{\textcolor{stewartcyan}{VÍ DỤ 10}}}\quad Tìm đạo hàm của mỗi hàm số sau:
\vspace{0.15cm}

\noindent
\textbf{(a)} $g(x) = 2^x$ \qquad\qquad\qquad \textbf{(b)} $h(x) = 5^{-x}$

\vspace{0.25cm}
\noindent\textbf{\textsf{\textcolor{stewartcyan}{LỜI GIẢI}}}
\vspace{0.1cm}

\noindent
\textbf{(a)} Ta sử dụng Công thức 5 với $b = 2$:
\[
g'(x) = \frac{d}{dx}(2^x) = 2^x \ln 2
\]
Kết quả này phù hợp với ước lượng
\[
\frac{d}{dx}(2^x) \approx (0{,}693)2^x
\]
mà chúng ta đã đưa ra trong \textbf{\textsf{Mục 3.1}} vì $\ln 2 \approx 0{,}693147$.

\end{minipage}
